\chapter{Implementation}
\label{ch:implementation}

% Structure mirrors the host lab's thesis format (3 sections):
%   4.1 Technologies
%   4.2 Software Architecture
%   4.3 UML Flow diagrams

This chapter documents how the design decisions of
Chapter~\ref{ch:requirements-design} were realised in working software.
Section~\ref{sec:impl-technologies} surveys the technology stack.
Section~\ref{sec:impl-architecture} describes the software architecture
of AIutami, with emphasis on the backend split into purpose-specific
Django apps. Section~\ref{sec:impl-uml} presents the UML diagrams that
capture the runtime behaviour of the most relevant subsystems
(turn-taking, audio forwarding, moderation pipeline, session
lifecycle).

\section{Technologies}
\label{sec:impl-technologies}
% TODO. Suggested coverage:
%   - Backend framework: Django + Daphne (ASGI), Channels for WebSockets
%   - Real-time audio: aiortc (WebRTC server-side), STUN/TURN considerations
%   - AI services: OpenAI Realtime (STT), OpenAI gpt-4o-mini (LLM),
%                  OpenAI gpt-4o-mini-tts (TTS)
%   - Persistence: PostgreSQL 16 (sessions, participants, votes, events),
%                  Redis 7 (turn state, ASR cache, moderation state)
%   - Frontend (separate repo): React + Vite + Tailwind, mobile-first PWA
%   - Deploy: Docker Compose, Nginx host-level, Let's Encrypt HTTPS
%   - Reporting: ReportLab (PDF), per-task report templates
% Justify each choice briefly (why Daphne over uvicorn, why aiortc over
% mediasoup, etc.).

\section{Software Architecture}
\label{sec:impl-architecture}
% TODO. Suggested coverage:
%   - High-level diagram (web/db/redis containers + WebRTC ports + nginx host)
%   - Django app split (apps/sessions, apps/turns, apps/asr, apps/tts,
%     apps/moderation, apps/webrtc, apps/reports, apps/tasks)
%   - Session lifecycle: LOBBY -> INDIVIDUAL_RANKING -> ACTIVE -> CONCLUSION -> CLOSED
%   - Task plugin system (registry, base.py, four implementations)
%   - WebSocket endpoints: /ws/sessions/, /ws/turns/, /ws/webrtc/
%   - Turn-taking state machine (IDLE / HUMAN_SPEAKING / AI_SPEAKING /
%     AI_INTRODUCING) and reservation window
%   - Audio forwarding hub (no mixing, per-peer outbound queue,
%     backpressure via drop chunk)
%   - ASR pipeline (gating, OpenAI Realtime WebSocket, transcript persistence)
%   - TTS pipeline (gpt-4o-mini-tts, resample 24k->48k)
%   - Moderation orchestrator (triggers, LLM, intervention decisions)
%   - No-moderator mode (control arm) — guards in coordinator only
%   - GDPR consent flow, GDPR-compliant data retention

\section{UML Flow Diagrams}
\label{sec:impl-uml}
% TODO. Suggested diagrams:
%   - Sequence diagram: human turn lifecycle (frontend -> turns WS -> ASR ->
%     moderation orchestrator -> optional TTS -> back to IDLE)
%   - State diagram: TurnManager states + reservation window
%   - Sequence diagram: AI moderator intervention triggered by NO_PUSH timer
%   - Sequence diagram: forced conclusion at TIMER_30 expiry
%   - Activity diagram: session lifecycle with task-specific branches
%     (individual ranking phase only for survival tasks)
